\section{Plano de ação, indicadores e gatilhos}
\begin{table}[H]\centering
\caption{Plano de ação resumido da Entrega XIII}
\small
\begin{tabularx}{\textwidth}{L{2.5cm}Y L{2.5cm}Y}
\toprule
Frente & Ações & Recursos & Indicadores\\\midrule
E1 -- praça piloto & Entrevistas, 10 testes sem ajuda, 2--3 clínicas, QR por parceiro e acompanhamento concierge. & Tempo da equipe; materiais; consultoria veterinária. & 15--20 entrevistas; 8/10 tarefas concluídas; 40--50 participantes convidados; 2 clínicas ativas.\\
E2 -- monetização & Testar preço, adesão ao PetCuida+ e Plano Parceiro; negociar patrocínio formal; limitar créditos. & Formulário, material comercial, apoio jurídico e financeiro. & Conversão paga; MRR recebido; cancelamento; margem por atendimento; custo do fundo.\\
Tecnologia e proteção & Corrigir controles de acesso, consentimento, logs, backup e testes; automatizar fluxos estáveis. & Desenvolvimento, revisão de segurança e infraestrutura. & 0 vulnerabilidades críticas abertas; 100\% dos fluxos críticos testados; custo/MAU.\\
Replicação & Documentar kit de praça e checklist de implantação. & Manual, contratos, treinamento e painel. & Tempo de implantação; custo por praça; qualidade e retenção local.\\
\bottomrule
\end{tabularx}\fonteautores
\end{table}

\subsection{Definições operacionais}
\begin{itemize}
\item \textbf{Ativação D7}: usuário completa perfil do animal e cadastra pelo menos um lembrete em até sete dias após o cadastro.
\item \textbf{Retenção D30}: usuário realiza ao menos uma ação de valor entre os dias 22 e 30 (consulta ao histórico, atualização de cuidado preventivo ou encaminhamento).
\item \textbf{CAC ativado}: gastos atribuíveis ao canal divididos pelo número de usuários que atingiram ativação D7 no mesmo período/coorte.
\item \textbf{Conversão paga}: assinantes pagantes ativos divididos por usuários ativos elegíveis; excluir testes gratuitos.
\item \textbf{Margem por atendimento}: comissão efetivamente recebida menos taxas, impostos, custos variáveis e fundo de créditos efetivamente liquidado.
\end{itemize}

\subsection{Gatilhos para avançar}
A equipe só deverá abrir uma segunda praça quando, por pelo menos dois ciclos mensais consecutivos: (i) houver oferta local contratada e disponibilidade confirmada; (ii) ativação D7 e retenção D30 forem medidas em coortes com denominadores explícitos; (iii) CAC observado for compatível com payback e CLV; (iv) margem por atendimento não for negativa ou o subsídio tiver fonte contratada; (v) houver receita recorrente recebida e não apenas intenção; e (vi) segurança, privacidade e suporte estiverem operacionais. As metas numéricas das versões anteriores (10 clínicas, 1.500 usuários, 5\% assinantes) permanecem referências aspiracionais, não critérios suficientes por si sós.